% Part: turing-machines
% Chapter: undecidability
% Section: unsolvability-decision-problem

\documentclass[../../../include/open-logic-section]{subfiles}

\begin{document}

\olfileid[pt]{tur}{und}{dec}
\olsection{O Problema da Decisão}

Dizemos que a lógica de primeira ordem é \emph{decidível} se, e
somente se, há um método efetivo para determinar se uma dada
!!{sentence} é válida ou não. Acontece que não existe tal método: o
problema de decidir a validade de sentenças de primeira ordem é
insolúvel.

A fim de estabelecer esse importante resultado negativo, provamos que
o problema da decisão não pode ser resolvido por uma máquina de
Turing. Isto é, mostramos que não existe máquina de Turing que, sempre
que iniciada em uma fita que contém uma !!{sentence} de primeira
ordem, eventualmente para e produz $1$ ou~$0$ conforme a
!!{sentence} seja válida ou não. Pela tese de Church-Turing, toda
função que é computável é Turing-computável. Logo, se essa ``função de
validade'' fosse de algum modo efetivamente computável, ela seria
Turing-computável. Se ela não é Turing-computável, então também não
pode ser efetivamente computável.

Nossa estratégia para provar que o problema da decisão é insolúvel é
reduzir o problema da parada a ele. Isso significa o seguinte:
provamos que a função~$h(e,w)$ que para com saída~$1$ se a máquina de
Turing descrita por~$e$ para na entrada~$w$ e produz~$0$ caso
contrário não é Turing-computável. Mostraremos que se houvesse uma
máquina de Turing que decide a validade de sentenças de primeira
ordem, então haveria também uma máquina de Turing que computa~$h$.
Como $h$ não pode ser computada por uma máquina de Turing, também não
pode haver uma máquina de Turing que decide a validade.

O primeiro passo nessa estratégia é mostrar que, para toda entrada~$w$
e máquina de Turing~$M$, podemos descrever efetivamente !!a{sentence}
$!T(M, w)$ representando o conjunto de instruções de~$M$ e a
entrada~$w$, e !!a{sentence}~$!E(M, w)$ exprimindo ``$M$ eventualmente
para'', tal que:
\begin{quote}
  $\Entails !T(M, w) \lif !E(M,w)$ se, e somente se, $M$ para na
  entrada~$w$.
\end{quote}
A maior parte de nossa prova consistirá em descrever essas sentenças
$!T(M, w)$ e~$!E(M, w)$ e em verificar que $!T(M, w) \lif !E(M, w)$
é válida se, e somente se, $M$~para na entrada~$w$.

\end{document}
